\documentclass[10pt]{beamer}
\usetheme{Madrid}
\usecolortheme{whale}

\usepackage[utf8]{inputenc}
\usepackage{booktabs}
\usepackage{tabularx}
\usepackage{graphicx}
\usepackage{listings}
\usepackage{tikz}

% Configuration for code snippet presentations if needed
\lstset{
    backgroundcolor=\color{gray!10},
    basicstyle=\ttfamily\scriptsize,
    breaklines=true,
    keywordstyle=\color{blue},
    commentstyle=\color{green!50!black},
    stringstyle=\color{orange}
}

% --- Presentation Info ---
\title[PitchVision]{PitchVision: Broadcast to 2D Mini-Map Projection}
\subtitle{Computer Vision Course Project}
\author[Peter Yaacoub]{Peter Yaacoub}
\institute[Computer Vision]{Project Type: Productivity}
\date{\today}

\begin{document}

% --- Slide 1: Title Slide ---
\begin{frame}
    \titlepage
\end{frame}

% --- Slide 2: Project Motivation & Objectives ---
\begin{frame}{Project Motivation \& Goal}
    \begin{block}{Problem Statement}
        Broadcast cameras provide excellent viewing perspective but fail to explicitly convey the full tactical layout and player configurations across a pitch.
    \end{block}
    \vspace{0.3cm}
    \textbf{Project Goal:} \\
    Build an experimental, unsupervised pipeline that takes standard broadcast frames, processes them, and projects player coordinates onto a standardized 2D tactical mini-map.
    \vspace{0.3cm}
    \begin{itemize}
        \item No camera calibration metrics provided beforehand.
        \item Operates without verified ground-truth annotations.
        \item Identifies explicit failure modes and practical limitations of unsupervised frameworks.
    \end{itemize}
\end{frame}

% --- Slide 3: Scope and Non-Goals ---
\begin{frame}{Scope \& Boundaries}
    To keep the pipeline strictly focused on core Computer Vision components:
    \vspace{0.4cm}
    \begin{columns}[t]
        \begin{column}{0.48\textwidth}
            \color{blue}\textbf{What PitchVision Does:}
            \begin{itemize}
                \color{black}
                \item Detects players and referee via YOLO.
                \item Clusters jerseys by color groups.
                \item Generates homography matrices from field line segmentations.
                \item Tracks player ids across consecutive frames.
            \end{itemize}
        \end{column}
        \begin{column}{0.48\textwidth}
            \color{red}\textbf{Explicit Non-Goals:}
            \begin{itemize}
                \color{black}
                \item No predictive modeling on future movements/passes.
                \item No multi-camera sync (single continuous shot focus).
                \item Does not try to create perfect reconstructions over entire matches.
            \end{itemize}
        \end{column}
    \end{columns}
\end{frame}

% --- Slide 4: Overall Pipeline Architecture ---
\begin{frame}{The Implemented Pipeline}
    The framework operates sequentially on raw video inputs:
    \vspace{0.2cm}
    \begin{enumerate}
        \item \textbf{Frame Extraction:} Sampled evenly spaced subsets (for debugging) or sequential intervals (for final analysis).
        \item \textbf{Object Detection:} Deep learning (YOLO) identifies players, ball, and referees.
        \item \textbf{Jersey-Color Clustering:} Masking out the grass backdrop and identifying team boundaries via K-Means.
        \item \textbf{Field Line Extraction:} Canny Edge + Hough Transform to determine layout indicators.
        \item \textbf{Approximate Projection:} Direct transformation matrix mapping frame boundaries onto a 2D plane.
        \item \textbf{Tracking:} Smooth frame-to-frame correlations with persistent track IDs.
    \end{enumerate}
\end{frame}

% --- Slide 5: Detection & Jersey-Color Clustering ---
\begin{frame}{Detection \& Team Clustering}
    \begin{itemize}
        \item \textbf{YoloDetector:} Crops bounding boxes around players and filters them based on expected bounding box area parameters.
        \item \textbf{Grass Masking:} Uses an HSV range filter (35--95 Hue) to isolate and mask out field green, protecting color integrity.
        \item \textbf{Upper Body Isolation:} Crops to the upper $62\%$ of the player box to minimize socks/boots interference.
        \item \textbf{Clustering Strategy:} Employs K-Means to identify primary team configurations while handling exceptions for referees and goalkeepers.
    \end{itemize}
\end{frame}

% --- Slide 6: Field Masking \& Homography Estimation ---
\begin{frame}{Pitch Line Extraction \& Homography}
    Establishing spatial coordinates without fixed calibrations:
    \begin{itemize}
        \item \textbf{Field Segments:} Extracted via morphology operators on green masks.
        \item \textbf{Line Detection:} Hough Transform (\texttt{HoughLinesP}) isolates white pitch markings.
        \item \textbf{Homography Estimation:} Map visual boundaries to standard $105\text{m} \times 68\text{m}$ definitions.
    \end{itemize}
    \vspace{0.1cm}
    \begin{figure}
        \centering
        \includegraphics[width=0.75\textwidth]{../../outputs/final/overlays/frame_000001.jpg}
        \caption{Visualizing extracted field lines and player detection overlays.}
    \end{figure}
\end{frame}

% --- Slide 7: Multiple Object Tracking System ---
\begin{frame}{The SimpleTracker Framework}
    \begin{block}{KISS Paradigm Implementation}
        Swapped complex Hungarian/Kalman setups for a light, deterministic, dependency-free tracker relying on \textbf{constant-velocity predictions} and \textbf{greedy assignment}.
    \end{block}
    \vspace{0.2cm}
    \textbf{Assignment Cost Metric:}
    \[ \text{Cost} = \text{Spatial Distance} + \text{Team Penalty} + \text{Color Penalty} \]
    \begin{itemize}
        \item \textbf{Velocity Vector:} Rolling momentum: $v_{new} = 0.6 \cdot v_{curr} + 0.4 \cdot v_{old}$.
        \item \textbf{Persistence:} Tracks persist up to $6$ missing frames to handle temporary occlusion.
    \end{itemize}
\end{frame}

% --- Slide 8: Data Splits and Evaluation Setup ---
\begin{frame}{Experimental Setup \& Datasets}
    The code supports subcommands: \texttt{extract-frames}, \texttt{run-checkpoint}, and \texttt{run}.
    \vspace{0.3cm}
    \begin{itemize}
        \item \textbf{50\% Milestone Subset:} 50 static frames for validating detection and clustering.
        \item \textbf{Final Target Segments:} Continuous video (10--15s) for evaluating panning artifacts and tracking consistency.
    \end{itemize}
    \vspace{0.2cm}
    \textbf{Metrics:} Team clustering F1-scores and relative projection stability.
\end{frame}

% --- Slide 9: Intermediate vs Final Milestones ---
\begin{frame}{Milestone Progression \& Visualization}
    \begin{columns}[t]
        \begin{column}{0.48\textwidth}
            \textbf{50\% Checkpoint Status:}
            \begin{itemize}
                \item Reliable player bounding boxes.
                \item Upper torso extraction resolved class mixing.
                \item Ball detection thresholds optimized.
            \end{itemize}
        \end{column}
        \begin{column}{0.48\textwidth}
            \textbf{Final Implementation Outputs:}
            \begin{itemize}
                \item Stitched side-by-side videos.
                \item Live updating mini-maps.
                \item Unified logging via JSON.
            \end{itemize}
        \end{column}
    \end{columns}
    \vspace{0.4cm}
    \begin{figure}
        \centering
        \includegraphics[width=0.85\textwidth]{../../outputs/final/minimap_frames/frame_000001.jpg}
        \caption{Side-by-side view: Broadcast overlay vs tactical tracking mini-map.}
    \end{figure}
\end{frame}

% --- Slide 10: System Limitations ---
\begin{frame}{Project Evaluation \& Limitations}
    \begin{block}{Observed Failure Modes}
        \begin{itemize}
            \item \textbf{Projection Jitter:} Slight changes in line detections skew the homography matrix calculation frame-to-frame.
            \item \textbf{Jersey Mixing:} Shadows, movement, or kits close to green create clustering edge cases.
        \end{itemize}
    \end{block}
\end{frame}

\end{document}